%%%%%%%%%%%%%%%%%%%%%%%%%%%%%%%%%%%%%%%%%%%%%%%%%%%%%%%%%%%%%%%%%%%%%%%%%%%%%%%%
%%%%%%%%%%%%%%%%%%%%%%%%%%%%%%%%%%%%%%%%%%%%%%%%%%%%%%%%%%%%%%%%%%%%%%%%%%%%%%%%
%%                                                                            %%
%% thesistemplate.tex version 4.01 (2023/09/21)                               %%
%% The LaTeX template file to be used with the aaltothesis.sty (version 4.00) %%
%% style file.                                                                %%
%% This package requires pdfx.sty v. 1.5.84 (2017/05/18) or newer.            %%
%%                                                                            %%
%% This is licensed under the terms of the MIT license below.                 %%
%%                                                                            %%
%% Written by Luis R.J. Costa.                                                %%
%% Currently developed at Teacher services, Learning Services of Aalto        %%
%% University by Luis R.J. Costa since May 2019.                              %%
%%                                                                            %%
%% Copyright 2017-2021 aaltothesis.cls by Luis R.J. Costa,                    %%
%% luis.costa@aalto.fi.                                                       %%
%% Copyright 2017-2018 Swedish translations in aaltothesis.cls by Elisabeth   %%
%% Nyberg and Henrik Wallén henrik.wallen@aalto.fi.                           %%
%% Finnish documentation in the template opinnatepohja.tex translated from    %%
%% the English template documentation.                                        %%
%% Copyright 2021 English template thesistemplate.tex by Luis R.J. Costa,     %%
%% Maurice Forget, Henrik Wallén.                                             %%
%% Copyright 2018-2022 Swedish template kandidatarbetsbotten.tex by Henrik    %%
%% Wallen.                                                                    %%
%%                                                                            %%
%% Permission is hereby granted, free of charge, to any person obtaining a    %%
%% copy of this software and associated documentation files (the "Software"), %%
%% to deal in the Software without restriction, including without limitation  %%
%% the rights to use, copy, modify, merge, publish, distribute, sublicense,   %%
%% and/or sell copies of the Software, and to permit persons to whom the      %%
%% Software is furnished to do so, subject to the following conditions:       %%
%% The above copyright notice and this permission notice shall be included in %%
%% all copies or substantial portions of the Software.                        %%
%% THE SOFTWARE IS PROVIDED "AS IS", WITHOUT WARRANTY OF ANY KIND, EXPRESS OR %%
%% IMPLIED, INCLUDING BUT NOT LIMITED TO THE WARRANTIES OF MERCHANTABILITY,   %%
%% FITNESS FOR A PARTICULAR PURPOSE AND NONINFRINGEMENT. IN NO EVENT SHALL    %%
%% THE AUTHORS OR COPYRIGHT HOLDERS BE LIABLE FOR ANY CLAIM, DAMAGES OR OTHER %%
%% LIABILITY, WHETHER IN AN ACTION OF CONTRACT, TORT OR OTHERWISE, ARISING    %%
%% FROM, OUT OF OR IN CONNECTION WITH THE SOFTWARE OR THE USE OR OTHER        %%
%% DEALINGS IN THE SOFTWARE.                                                  %%
%%                                                                            %%
%%                                                                            %%
%%%%%%%%%%%%%%%%%%%%%%%%%%%%%%%%%%%%%%%%%%%%%%%%%%%%%%%%%%%%%%%%%%%%%%%%%%%%%%%%
%%                                                                            %%
%%                                                                            %%
%% An example for writting your thesis using LaTeX                            %%
%% Original version and development work by Luis Costa, changes to the text   %% 
%% in the Finnish template by Perttu Puska.                                   %%
%% Support for Swedish added 15092014                                         %%
%% PDF/A-b support added on 15092017                                          %%
%% PDF/A-2 support added on 24042018                                          %%
%% Layout design and typesettin changed 15072021                              %%
%%                                                                            %%
%% This example consists of the files                                         %%
%%       thesistemplate.tex (version 4.00) (for text in English)              %%
%%       opinnaytepohja.tex (version 4.00) (for text in Finnish)              %%
%%       kandidatarbetsbotten.tex (version 1.10) (for text in Swedish)        %%
%%       aaltothesis.cls                                                      %%
%%       linediagram.pdf (graphics file)                                      %%
%%       curves.pdf      (graphics file)                                      %%
%%       ledspole.jpg    (graphics file)                                      %%
%%                                                                            %%
%%                                                                            %%
%% Typeset in Linux with                                                      %%
%% pdflatex: (recommended method)                                             %%
%%             $ pdflatex thesistemplate                                      %%
%%             $ pdflatex thesistemplate                                      %%
%%                                                                            %%
%%   The result is the file thesistemplate.pdf that is PDF/A compliant, if    %%
%%   you have chosen the proper \documenclass options (see comments below)    %%
%%   and your included graphics files have no problems.                       %%
%%                                                                            %%
%%                                                                            %%
%% Explanatory comments in this example begin with the characters %%, and     %%
%% changes that the user can make with the character %                        %%
%%                                                                            %%
%%%%%%%%%%%%%%%%%%%%%%%%%%%%%%%%%%%%%%%%%%%%%%%%%%%%%%%%%%%%%%%%%%%%%%%%%%%%%%%%
%%%%%%%%%%%%%%%%%%%%%%%%%%%%%%%%%%%%%%%%%%%%%%%%%%%%%%%%%%%%%%%%%%%%%%%%%%%%%%%%
%%
%% WHAT is PDF/A
%%
%% PDF/A is the ISO-standardized version of the pdf. The standard's goal is to
%% ensure that he file is reproducable even after a long time. PDF/A differs
%% from pdf in that it allows only those pdf features that support long-term
%% archiving of a file. For example, PDF/A requires that all used fonts are
%% embedded in the file, whereas a normal pdf can contain only a link to the
%% fonts in the system of the reader of the file. PDF/A also requires, among
%% other things, data on colour definition and the encryption used.
%% Currently three PDF/A standards exist:
%% PDF/A-1: based on PDF 1.4, standard ISO19005-1, published in 2005.
%%          Includes all the requirements essential for long-term archiving.
%% PDF/A-2: based on PDF 1.7, standard ISO19005-2, published in 2011.
%%          In addition to the above, it supports embedding of OpenType fonts,
%%          transparency in the colour definition and digital signatures.
%% PDF/A-3: based on PDF 1.7, standard ISO19005-3, published in 2012.
%%          Differs from the above only in that it allows embedding of files in
%%          any format (e.g., xml, csv, cad, spreadsheet or wordprocessing
%%          formats) into the pdf file.
%% PDF/A-4: based on PDF 2.0, standard ISO19005-4, published in November 2020.
%%
%% PDF/A-1 files are not necessarily PDF/A-2 -compatible and PDF/A-2 are not
%% necessarily PDF/A-1 -compatible.
%% Standards PDF/A-1, PDF/A-2 and PDF/A-3 have two levels:
%% b: (basic) requires that the visual appearance of the document is reliably
%%    reproduceable.
%% a (accessible) in addition to the b-level requirements, specifies how
%%   accessible the pdf file is to assistive software, say, for the physically
%%   impaired.
%% The PDF/A-4 standard does not have additional levels like in the earlier
%% standards.
%% For more details on PDF/A, see, e.g., 
%% https://www.loc.gov/preservation/digital/formats/fdd/fdd000318.shtml or
%% https://www.pdfa.org/resource/iso-19005-pdfa/
%%
%%
%% WHICH PDF/A standard should my thesis conform to?
%%
%% Either to the PDF/A-1b or the PDF/A-2b standard. If all the figures and
%% graphs used in thesis work do not require transparency features, use either
%% PDF/A-1b or PFDF/A-2b. If you have figures with transparency
%% characteristics, use the PDF/A-2b standard. However, drawing applications
%% often use the transparency parameter, setting it to zero, to specify opacity
%% and get the basic 2-D visualisation. As a result, validation of PDF/A-1b
%% will fail. Use PDF/A-2b if PDF/A-1b validation fails.
%% Do not use the PDF/A-3b standard for your thesis.
%% The font to be used are specified in this templatenand they should not be
%% changed. In addition to not adhering to Aalto's visual guidelines, you may
%% have difficulties in producing a PDF/A-compliant PDF.
%%
%%
%% Validate your PDF/A file at https://www.pdf-online.com/osa/validate.aspx
%%
%%
%% WHAT graphics format can I use to produce my PDF/A compliant file?
%%
%% When using pdflatex to compile your work, favour the use of pdf, but you can
%% use the jpg or png format especially for photographs. You will have PDF/A 
%% compliance problems with figures in pdf if the fonts are not embedded in the
%% pdf file.
%% If you choose to use latex to compile your work, the only acceptable file
%% format for your figure is eps. DO NOT use the ps format for your figures.

%% USE one of the following three \documentclass set-ups:
%% * the first when using pdflatex to directly typeset your document in the
%%   chosen pdf/a format for online publishing (centred page layout),
%% * the second for one-sided printing your thesis with the layout (wide left 
%%   margin), or
%% * the third for two-sided printing.
%%
\documentclass[english, 12pt, a4paper, biz, utf8, a-2b, online]{aaltothesis}
%\documentclass[english, 12pt, a4paper, elec, utf8, a-2b, print]{aaltothesis}
%\documentclass[english, 12pt, a4paper, elec, utf8, a-2b, print, twoside]{aaltothesis}

\renewcommand{\AaltoLogoSmall}[3]{%
    \includegraphics[width=0.3\textwidth]{./figures/logos/EFREI.png} % Replace with your logo path
}

%% Use the following options in the \documentclass macro above:
%% your school: arts, biz, chem, elec, eng, sci
%% the character encoding scheme used by your editor: utf8, latin1
%% thesis language: english, finnish, swedish
%% make an archiveable PDF/A-1b or PDF/A-2b compliant file: a-1b, a-2b
%%                    (with pdflatex, a normal pdf containing metadata is
%%                     produced without the a-*b option)
%% typset for online document or print on paper: online, print
%%        online: typeset in symmetric layout and blue hypertext for online
%%                publishing
%%        print: typeset in a symmetric layout and black hypertext for printing
%%               on paper
%%          two-side printing: twoside (default is one-sided printing)
%%               typeset in a wide margin on the binding side of the page and
%%               black hypertext. Use with print only.
%%

%% Use one of these if you write in Finnish (or use the Finnish template
%% opinnaytepohja.tex)
%\documentclass[finnish, 12pt, a4paper, elec, utf8, a-1b, online]{aaltothesis}
%\documentclass[finnish, 12pt, a4paper, elec, utf8, a-1b, print]{aaltothesis}
%\documentclass[finnish, 12pt, a4paper, elec, utf8, a-1b, print, twoside]{aaltothesis}

%% Use one of these if you write in Swedish (or use the Swedish template
%% kandidatarbetsbotten.tex)
%\documentclass[swedish, 12pt, a4paper, elec, utf8, a-2b, online]{aaltothesis}
%\documentclass[swedish, 12pt, a4paper, elec, utf8, a-2b]{aaltothesis}
%\documentclass[swedish, 12pt, a4paper, elec, dvips, online]{aaltothesis}

%% FOR USERS OF AMS PACKAGES:
%% * newtxmath used in this template loads amsmath, so
%%   you needn't load it. If you want to use options in amsmath, load it here, 
%%   before \setupthesisfonts below to pass the options to amsmath.
%% * If you want to use amsthm, load it here before \setupthesisfonts to avoid
%%   a clash with newtxmath.
%% * If using amsmath with options and you want to use amsthm, load amsthms
%%   after amsmath, as described in the amsthm documentation.
%% * Don't use amsbsym or amsfonts. The symbols [and macros] there are defined in
%%   newtxmath and so clash if used.
%\usepackage[options]{amsmath}
%\usepackage{amsthm}
\usepackage{url}
\usepackage[utf8]{inputenc}
\usepackage[T1]{fontenc}
\usepackage{listings}
\usepackage{xcolor}
\usepackage{tcolorbox}
\tcbuselibrary{listings}

% Define custom colours
\definecolor{vscBackground}{rgb}{0.12, 0.12, 0.12} % Softer dark background
\definecolor{vscText}{rgb}{0.96, 0.96, 0.96}       % Light grey text
\definecolor{vscKeywordBlue}{rgb}{0.34, 0.61, 0.84}% Blue for definitions (RGB: 86,156,214)
\definecolor{vscControlPink}{rgb}{0.77, 0.52, 0.75}% Pink for control flow (RGB: 197,134,192)
\definecolor{vscOptionGreen}{rgb}{0.3059, 0.7882, 0.6902}
\definecolor{vscCommentGray}{rgb}{0.5, 0.5, 0.5}   % Comments in grey
\definecolor{vscNumberGray}{rgb}{0.8, 0.8, 0.8}    % Line numbers in light grey
\definecolor{vscStringOrange}{rgb}{0.95, 0.65, 0.2}% Strings in orange (optional)

% Listings style configuration
\lstdefinelanguage{Coq}{
  morekeywords=[1]{Module, Definition, Some, None, End}, % Group 1: Blue
  morekeywords=[2]{match, with, if, then, else, end},    % Group 2: Pink
  morekeywords=[3]{option, Type}, % Group 3: Green
  sensitive=true,
  morecomment=[l]{(*}, 
  morecomment=[s]{(*}{*)},
  morestring=[b]", % Strings
}

\lstset{
  language=Coq,
  basicstyle=\ttfamily\footnotesize\color{vscText},  % Default text style
  keywordstyle=[1]\color{vscKeywordBlue},   % Group 1: Blue keywords
  keywordstyle=[2]\color{vscControlPink},   % Group 2: Pink keywords
  keywordstyle=[3]\color{vscOptionGreen},
  commentstyle=\itshape\color{vscCommentGray},       % Comments in italic grey
  stringstyle=\color{vscStringOrange},               % Strings in orange
  backgroundcolor=\color{vscBackground},             % Dark background
  frame=none,                                        % Remove visible border
  framerule=0pt,                                     % No rule thickness
  rulecolor=\color{vscBackground},                   % Match frame color to background
  breaklines=true,                                   % Break long lines
  tabsize=2,
  showspaces=false,
  showstringspaces=false,
  aboveskip=10pt,
  belowskip=10pt,
  escapeinside={(*@}{@*)},
}

\renewcommand{\lstlistingname}{Figure}

%% DO NOT MOVE OR REMOVE \setupthesisfonts
\setupthesisfonts

%%
%% Add here the packges you need
%%
\usepackage{graphicx}


%% For tables that span multiple pages; used to split a paraphrasing example in
%% the appendix. If you don't need it, remove it.
\usepackage{longtable}

%% A package for generating Creative Commons copyright terms. If you don't use
%% the CC copyright terms, remove it, since otherwise undesired information may
%% be added to this document's metadata.
\usepackage[type={CC}, modifier={by-nc-sa}, version={4.0}]{doclicense}
%% Find below three examples for typesetting the CC license notice.


%% Edit to conform to your degree programme
%% Capitalise the words in the name of the degree programme: it's a name
\degreeprogram{Cybersecurity and Management}
%%

%% Your major
%%
\major{Cybersecurity}
%%

%% Choose one of the three below
%%
%\univdegree{BSc}
\univdegree{MSc}
%\univdegree{Lic}
%%

%% Your name (self explanatory...)
%%
\thesisauthor{Antoine James}
%%

%% Your thesis title and possible subtitle comes here and possibly, again,
%% together with the Finnish or Swedish abstract. Do not hyphenate the title
%% (and subtitle), and avoid writing too long a title. Should LaTeX typeset a
%% long title (and/or subtitle) unsatisfactorily, you might have to force a
%% linebreak using the \\ control characters. In this case...
%% * Remember, the title should not be hyphenated!
%% * A possible 'and' in the title should not be the last word in the line; it
%%   begins the next line.
%% * Specify the title (and/or subtitle) again without the linebreak characters
%%   in the optional argument in box brackets. This is done because the title
%%   is part of the metadata in the pdf/a file, and the metadata cannot contain
%%   linebreaks.
%%
\thesistitle{Formal verification of blockchain protocols for correctness and reliability}
%\thesistitle[Title of the thesis]{Title of\\ the thesis}
%%
%% Either remove or leave \thesissubtitle{} empty if you don't use it
%%
%\thesissubtitle[Subtitle of the thesis]{Subtitle of\\ the thesis}
%\thesissubtitle{}

%%
\place{Espoo}
%%

%% The date for the bachelor's thesis is the day it is presented
%%
\date{21 March 2026}
%%

%% Thesis supervisor
%% Note the "\" character in the title after the period and before the space
%% and the following character string.
%% This is because the period is not the end of a sentence after which a
%% slightly longer space follows, but what is desired is a regular interword
%% space.
%%
\supervisor{Prof.\ TBD.}
%%

%% Advisor(s)---two at the most---of the thesis. Check with your supervisor how
%% many official advisors you can have.
%%
\advisor{Dr Guillaume Claret, Danny Willems}
%%

%% If you do your thesis work in a company of other institute, give the name of
%% the company or instution here. Otherwise, leave the macro empty, comment it
%% out, or remove it. This will remove this field from the abstract page.
%%
\collaborativepartner{Formal Land}
%%

%% Aaltologo: syntax:
%% \uselogo{?|!|''}
%% The logo language is set to be the same as the thesis language.
%%
%\uselogo{?}
%\uselogo{!}
\uselogo{''}
%%

%%%%%%%%%%%%%%%%%%               COPYRIGHT TEXT               %%%%%%%%%%%%%%%%%%
%%%%%%%%%%%%%%%%%%%%%%%%%%%%%%%%%%%%%%%%%%%%%%%%%%%%%%%%%%%%%%%%%%%%%%%%%%%%%%%%

%% Copyright of a work is with the creator/author of the work regardless of
%% whether the copyright mark is explicitly in the work or not. You may, if you
%% wish---we encourage you to do so---publish your work under a Creative
%% Commons license (see creativecommons.org), in which case the license text
%% must be visible in the work. Write here the copyright text you want using the
%% macro \copyrighttext, which writes the text into the metadata of the pdf file
%% as well.
%%
%% Syntax:
%% \copyrigthtext{metadata text}{text visible on the page}
%%
%% CHOOSE ONE OF THE COPYRIGHT NOTICE STYLES BELOW.
%% IF USING THE CC TERMS, CHOOSE THE LICENSE YOU WANT TO USE.
%% The different CC licenses are listed at 
%% https://creativecommons.org/about/cclicenses/.
%% If you use the icons from the dolicense.sty package, add the package above
%% (\usepackage{dolicense}).
%% IMPORTANT NOTE!! Manually write the CC text in the \copyrighttext metadata
%% text field.
%%
%% NOTE: In the macros below, the text written in the metadata must have a
%% \noexpand macro before the \copyright special character. When not in pdf/a
%% mode (i.e. a-1b or a-2b are not specified in \documentclass), two \noexpands
%% are required in the metadata text to correctly render the copyright mark in
%% the pdf metadata. In pdf/a mode one \noexpand suffices.
%%
%% EXAMPLE OF PLAIN COPYRIGHT TEXT
%% The macros \copyright and \year below must be separated by the \ character 
%% (space chacter) from the text that follows. The macros in the argument of the
%% \copyrighttext macro automatically insert the year and the author's name.
%% (Note! \ThesisAuthor is an internal macro of the aaltothesis.cls class file).
%%
%\copyrighttext{Copyright \noexpand\textcopyright\ \number\year\ \ThesisAuthor}
%{Copyright \textcopyright{} \number\year{} \ThesisAuthor}
%%
%% Of course, the same text could have simply been written as
%% \copyrighttext{Copyright \noexpand\copyright\ 2018 Eddie Engineer}
%% {Copyright \copyright{} 2022 Eddie Engineer}
%%
%% EXAMPLES OF CC LICENSE: different ways to display the same license
%% 1. A simple Creative Commons license text with a link to the copyright notice:
%\copyrighttext{\noexpand\textcopyright\ \number\year. This work is 
%	licensed under a CC BY-NC-SA 4.0 license.}{\textcopyright{} 
%	\number\year. This work is licensed under a 
%	\href{https://creativecommons.org/licenses/by-nc-nd/4.0/}{CC BY-NC-SA 4.0} 
%	license.}
%
%% To get the URL of the license of your choice, go to 
%% https://creativecommons.org/about/cclicenses/, click on the chosen license
%% you want to use, and copy-and-paste the URL in the macro \href above.
%%
%% 2. A short Creative Commons license text containing the respective CC icons
%% (requires the package dolicense.sty to be added in the preamble as done
%% above) and a link to the corresponding Creative Commons license webpage (see
%% the dolicense package documentation for other license icons):
%\copyrighttext{\noexpand\textcopyright\ \number\year. This work is licensed
%	under a CC BY-NC-SA 4.0 license.}{
%	\parbox{95mm}{\noindent\textcopyright\ \number\year. \doclicenseText} 
%	\hspace{1em}\parbox{35mm}{\doclicenseImage}
%}
%%
%% 3. An expanded Creative Commons license text containing the respective CC
%% icons text and as generated by the dolicense.sty package (the license is set
%% via package options in \usepackage[options]{dolicense} above; see the
%% dolicense package documentation for other license texts and icons):
\copyrighttext{\noexpand\textcopyright\ \number\year. This work is 
	licensed under a Creative Commons "Attribution-NonCommercial-ShareAlike 4.0 
	International" (BY-NC-SA 4.0) license.}{\noindent\textcopyright\ \number
	\year \ \doclicenseThis}
%%%%%%%%%%%%%%%%%%%%%%%%%%%%%%%%%%%%%%%%%%%%%%%%%%%%%%%%%%%%%%%%%%%%%%%%%%%%%%%%


%% The English abstract:
%% All the details (name, title, etc.) on the abstract page appear as specified
%% above.
%% Thesis keywords:
%% Note! The keywords are separated using the \spc macro
%%
\keywords{For keywords choose\spc concepts that are\spc central to your\spc 
thesis}
%%

%% The abstract text. This text in one paragraph is included in the metadata of
%% the pdf file as well as the abstract page. To have paragraphs in your
%% abstract rewrite it in the abstarct environment as described below.
%%
\thesisabstract{%
The abstract is a short description of the essential contents of the thesis
usually in one paragraph: what was studied and how and what were the main
findings. For a Finnish thesis, the abstract should be written in both Finnish
and English; for a Swedish thesis, in Swedish and English. The abstracts for
English theses written by Finnish or Swedish speakers should be written in
English and either in Finnish or in Swedish, depending on the student’s language
of basic education. Students educated in languages other than Finnish or Swedish
write the abstract only in English. Students may include a second or third
abstract in their native language, if they wish. 
The abstract text of this thesis is written on the readable abstract page as
well as into the pdf file's metadata via the thesisabstract macro (see the 
comment in the TeX file). Write here the text that goes into the metadata. The 
metadata cannot contain special characters, linebreak or paragraph break 
characters, so these must not be used here. If your abstract does not contain 
special characters and it does not require paragraphs, you may take advantage of
the abstracttext macro (see the comment in the TeX file below). Otherwise, the 
metadata abstract text must be identical to the text on the abstract page.
}

%% You can prevent LaTeX from writing into the xmpdata file (it contains all the 
%% metadata to be written into the pdf file) by setting the writexmpdata switch
%% to 'false'. This allows you to write the metadata in the correct format
%% directly into the file thesistemplate.xmpdata.
%\setboolean{writexmpdatafile}{false}


%% All that is printed on paper starts here
%%
\begin{document}

%% Create the coverpage
%%
\makecoverpage

%% Typeset the copyright text.
%% If you wish, you may leave out the copyright text from the human-readable
%% page of the pdf file. This may seem like a attractive idea for the printed
%% document especially if "Copyright (c) yyyy Eddie Engineer" is the only text
%% on the page. However, the recommendation is to print this copyright text.
%%
\makecopyrightpage

\clearpage
%% Note that when writing your thesis in English, place the English abstract
%% first followed by the possible Finnish or Swedish abstract.

%% Abstract text
%% All the details (name, title, etc.) on the abstract page appear as specified
%% above. Add your abstarct text with paragraphs here to have paragraphs in the
%% visible abstract page. Nonetheless, write the abstarct text without
%% paragraphs in the macro \thesismacro so that it is added to the metadata as
%% well.
%%
\begin{abstractpage}[english]
  The abstract is a short description of the essential contents of the thesis,
  usually in one paragraph: what was studied and how and what were the main
  findings.

  For a Finnish thesis, the abstract should be written in both Finnish and
  English; for a Swedish thesis, in Swedish and English. The abstracts for
  English theses written by Finnish or Swedish speakers should be written in
  English and either in Finnish or in Swedish, depending on the student’s
  language of basic education. Students educated in languages other than Finnish
  or Swedish write the abstract only in English. Students may include a second
  or third abstract in their native language, if they wish.

  The abstract text of this thesis is written on the readable abstract page as
  well as into the pdf file's metadata via the \verb+\thesisabstract+ macro
  (see comment in this \TeX{} file above). Write here the text that goes onto
  the readable abstract page. You can have special characters, linebreaks, and
  paragraphs here. Otherwise, this abstract text must be identical to the
  metadata abstract text.
  
  If your abstract does not contain special characters and it does not require
  paragraphs, you may take advantage of the \verb+\abstracttext+ macro (see the
  comment in this \TeX{} file below).
\end{abstractpage}

%% The text in the \thesisabstract macro is stored in the macro \abstractext, so
%% you can use the text metadata abstract directly as follows:
%%
%\begin{abstractpage}[english]
%	\abstracttext{}
%\end{abstractpage}

%% Force a new page so that the possible Finnish or Swedish abstract does not
%% begin on the same page
%%

\dothesispagenumbering{}

%% Preface
%%
%% This section is optional. Remove it if you do not want a preface.
\mysection{Preface}
%\mysection{Esipuhe}
I want to thank Professor Pirjo Professor and my instructors Dr Alan Advisor and
Ms Elsa Expert for their guidance.

I also want to thank my partner for keeping me sane and alive.

\vspace{5cm}
Otaniemi, 9 February 2023\\

\vspace{5mm}
{\hfill Eddie E.\ Engineer \hspace{1cm}}

%% Force a new page after the preface
%%
\newpage


%% Table of contents. 
%%
\thesistableofcontents


%% Symbols and abbreviations
\mysection{Symbols and abbreviations}

%% MODIFY THESE SYMBOLS %%

\subsection*{Symbols}

\begin{tabular}{ll}
$\mathbf{B}$  & magnetic flux density  \\
$c$              & speed of light in vacuum $\approx 3\times10^8$ [m/s]\\
$\omega_{\mathrm{D}}$    & Debye frequency \\
$\omega_{\mathrm{latt}}$ & average phonon frequency of lattice \\
$\uparrow$       & electron spin direction up\\
$\downarrow$     & electron spin direction down
\end{tabular}

\subsection*{Operators}

\begin{tabular}{ll}
$\nabla \times \mathbf{A}$              & curl of vectorin $\mathbf{A}$\\
$\displaystyle\frac{\mbox{d}}{\mbox{d} t}$ & derivative with respect to 
variable $t$\\[3mm]
$\displaystyle\frac{\partial}{\partial t}$  & partial derivative with respect 
to variable $t$ \\[3mm]
$\sum_i $                       & sum over index $i$\\
$\mathbf{A} \cdot \mathbf{B}$    & dot product of vectors $\mathbf{A}$ and 
$\mathbf{B}$
\end{tabular}

\subsection*{Abbreviations}

\begin{tabular}{ll}
AC         & alternating current \\
APLAC      & an object-oriented analog circuit simulator and design tool \\
           & (originally Analysis Program for Linear Active Circuits) \\
BCS        & Bardeen-Cooper-Schrieffer \\ %% dash between the names
DC         & direct current \\
TEM        & transverse eletromagnetic
\end{tabular}


%% \clearpage is similar to \newpage, but it also flushes the floats (figures
%% and tables).
%%
\cleardoublepage

%% Text body begins. Note that since the text body is mostly in Finnish the
%% majority of comments are also in Finnish after this point. There is no point
%% in explaining Finnish-language specific thesis conventions in English.
%% This text will be translated to English soon.
%%
\section{Introduction to Sui and the Layer 1 Blockchain Landscape}

\subsection{Definition of a Layer 1 Blockchain}
Layer 1 blockchains form the foundation of decentralised networks, acting as the primary infrastructure that enables secure, transparent, and immutable transactions without relying on intermediaries. These systems independently process and validate transactions, maintain the consistency of the network’s state, and support the execution of smart contracts directly on-chain. Unlike Layer 2 solutions, which function as off-chain enhancements designed to improve scalability or reduce transaction costs, Layer 1 blockchains address these challenges within their core architecture.

Notable examples of Layer 1 blockchains include Bitcoin, Ethereum, and Sui, each of which employs a distinct approach to balancing the trade-offs between scalability, decentralisation, and programmability. These trade-offs, often referred to as the \emph{blockchain trilemma}~\cite{werth2023review}, have shaped the design and evolution of Layer 1 protocols, with each attempting to overcome limitations while catering to specific use cases and user requirements.

\subsection{Overview of Sui as a Layer 1 Blockchain}
\emph{Sui}, designed by Mysten Labs, is a cutting-edge Layer 1 blockchain that redefines transaction processing through its \emph{object-centric data model}~\cite{suiobjects}. 

Unlike account-based models, such as those used by Ethereum, or UTXO-based systems like Bitcoin, Sui represents all on-chain entities—whether smart contracts, tokens, or digital assets—as independent \emph{objects}. Each object possesses a unique identifier and is directly associated with an account, eliminating the need for global state synchronisation. This allows Sui to focus on \emph{localised} object interactions.

Sui's object-centric model also facilitates \emph{parallel transaction execution}, wherein independent transactions can be processed simultaneously, greatly improving scalability. By departing from traditional paradigms, Sui achieves significant performance enhancements and opens up new possibilities for Web3 applications.

\subsection{The Move Programming Language}
The programmability of Sui is powered by the \texttt{Move} programming language, originally developed by Meta for the Libra blockchain. Move was designed with a strong emphasis on \emph{safety}, \emph{resource management}, and \emph{modularity}, making it an ideal choice for blockchain development.

Move's design is heavily influenced by the Rust programming language, renowned for its performance and memory safety. This foundation enables Move to inherit key attributes critical for blockchain applications. Specifically, Move's \emph{resource-oriented programming paradigm} is inspired by Rust's \emph{ownership model}~\cite{rust_ownership}, ensuring that digital assets are managed securely and efficiently. Moreover, Rust's emphasis on low-level performance contributes to Move's ability to execute smart contracts at scale, complementing Sui's object-centric architecture.

One of Move's key innovations is its treatment of \emph{resources as first-class citizens}, ensuring that digital assets cannot be duplicated or discarded inadvertently. In programming, treating something as a \emph{first-class citizen} means it has full capabilities to be created, passed to functions, returned from functions, and assigned to variables—just like standard data types such as integers or strings. This ensures that digital assets in Move are handled with strict safeguards, preventing common vulnerabilities such as double-spending or resource loss. Furthermore, Move's \emph{expressive type system} allows developers to encode and verify key properties at compile time, reducing runtime errors and vulnerabilities.

\subsection{Key Innovations in Sui}
Sui introduces several novel features that distinguish it from other Layer 1 blockchains:
\begin{itemize}
    \item \textbf{Object-Centric Data Model}: Simplifies transaction processing by representing on-chain entities as independent objects, enabling parallelism and scalability.
    \item \textbf{Parallel Transaction Execution}: Processes transactions concurrently to achieve higher throughput compared to sequential execution models.
    \item \textbf{Gas Pricing Mechanisms}: Employs dynamic gas pricing to balance user costs and validator incentives, creating a sustainable economic model.~\cite{sui_fees}
\end{itemize}

These innovations position Sui as a leader among next-generation blockchains, offering solutions to the longstanding challenges of scalability and efficiency.

\subsection{The Case for Formal Verification}
While Sui's innovations address many challenges in blockchain design, ensuring the security and correctness of the system remains paramount. Blockchain protocols and smart contracts are prone to errors and exploits, as evidenced by numerous high-profile security breaches in the decentralised finance (DeFi) sector.~\cite{phalcon_explorer}~\cite{rekt}

Formal verification provides a rigorous mathematical framework for proving the correctness of code. In the case of Sui, our work focused on verifying the \emph{type checker}, an essential component of the \texttt{Move} language, which ensures that programs adhere to their specified types and constraints. By leveraging tools such as \texttt{Coq}, the verification of the type checker guarantees critical properties, such as type safety and adherence to resource rules, under all conditions. Translating Rust code into formal specifications enables the creation of proofs that validate these properties, nullifying the risk of vulnerabilities.

This section establishes the groundwork for understanding Sui's architectural innovations and highlights the importance of formal verification in ensuring robust and secure blockchain protocols.

\begin{tcolorbox}[
  colback=white,     % Background color inside the box (set to white to match the document)
  colframe=white,    % Frame color (set to white to remove it)
  boxrule=0pt,       % No border
  sharp corners,     % Remove rounded corners
  boxsep=0pt,        % Remove padding
  left=0pt,          % Remove left margin
  right=0pt,         % Remove right margin
  top=0pt,           % Remove top margin
  bottom=0pt         % Remove bottom margin
]
\begin{lstlisting}[caption={Definition of the Inductive Type \texttt{t}},captionpos=b]
Module IntegerValue.
  Inductive (*@\textcolor{vscOptionGreen}{t}@*) : Type :=
    | U8 : Z -> t
    | U16 : Z -> t
    | U32 : Z -> t
    | U64 : Z -> t
    | U128 : Z -> t
    | U256 : Z -> t
    .
End IntegerValue.
\end{lstlisting}
\end{tcolorbox}

\textbf{Explanation}: The \texttt{Inductive t} defines a type \texttt{t} that supports multiple integer bit widths, such as \texttt{U8}, \texttt{U16}, \texttt{U32}, etc. Each constructor associates a bit length with a value of type \texttt{Z} (unbounded integers). This structure ensures that integers of varying sizes can be represented in the same type.

\begin{tcolorbox}[
  colback=white,     % Background color inside the box (set to white to match the document)
  colframe=white,    % Frame color (set to white to remove it)
  boxrule=0pt,       % No border
  sharp corners,     % Remove rounded corners
  boxsep=0pt,        % Remove padding
  left=0pt,          % Remove left margin
  right=0pt,         % Remove right margin
  top=0pt,           % Remove top margin
  bottom=0pt         % Remove bottom margin
]
\begin{lstlisting}[caption={Checked Addition \texttt{checked\_add}},captionpos=b]
Definition checked_add (a : ValueImpl.t) (b : ValueImpl.t) : option ValueImpl.t :=
  match a, b with
  | ValueImpl.U8 l, ValueImpl.U8 r =>
      if l + r <=? 2^8 - 1
      then Some (ValueImpl.U8 (l + r))
      else None
  | ValueImpl.U16 l, ValueImpl.U16 r =>
      if l + r <=? 2^16 - 1
      then Some (ValueImpl.U16 (l + r))
      else None
  | ValueImpl.U32 l, ValueImpl.U32 r =>
      if l + r <=? 2^32 - 1
      then Some (ValueImpl.U32 (l + r))
      else None
  | ValueImpl.U64 l, ValueImpl.U64 r =>
      if l + r <=? 2^64 - 1
      then Some (ValueImpl.U64 (l + r))
      else None
  | ValueImpl.U128 l, ValueImpl.U128 r =>
      if l + r <=? 2^128 - 1
      then Some (ValueImpl.U128 (l + r))
      else None
  | ValueImpl.U256 l, ValueImpl.U256 r =>
      if l + r <=? 2^256 - 1
      then Some (ValueImpl.U256 (l + r))
      else None
  | _, _ => None
  end.
\end{lstlisting}
\end{tcolorbox}
\clearpage
\textbf{Explanation}: The \texttt{checked\_add} function adds two integers of the same type (e.g., \texttt{U8}, \texttt{U16}). If the result exceeds the maximum value for that type, it returns \texttt{None}, preventing overflow. Otherwise, it wraps the result in a \texttt{Some} constructor.

\begin{tcolorbox}[
  colback=white,     % Background color inside the box (set to white to match the document)
  colframe=white,    % Frame color (set to white to remove it)
  boxrule=0pt,       % No border
  sharp corners,     % Remove rounded corners
  boxsep=0pt,        % Remove padding
  left=0pt,          % Remove left margin
  right=0pt,         % Remove right margin
  top=0pt,           % Remove top margin
  bottom=0pt         % Remove bottom margin
]
\begin{lstlisting}[caption={Checked Left Shift \texttt{shl\_checked}},captionpos=b]
Definition shl_checked (self : IntegerValue.t) (other : IntegerValue.t) : PartialVMResult.t IntegerValue.t :=
  match self, other with
  | IntegerValue.U8 l, IntegerValue.U8 r =>
      if r <? 8
      then Result.Ok (IntegerValue.U8 (Z.shiftl l r))
      else Result.Err (PartialVMError.new StatusCode.ARITHMETIC_ERROR)
  | IntegerValue.U16 l, IntegerValue.U16 r =>
      if r <? 16
      then Result.Ok (IntegerValue.U16 (Z.shiftl l r))
      else Result.Err (PartialVMError.new StatusCode.ARITHMETIC_ERROR)
  | IntegerValue.U32 l, IntegerValue.U32 r =>
      if r <? 32
      then Result.Ok (IntegerValue.U32 (Z.shiftl l r))
      else Result.Err (PartialVMError.new StatusCode.ARITHMETIC_ERROR)
  | IntegerValue.U64 l, IntegerValue.U64 r =>
      if r <? 64
      then Result.Ok (IntegerValue.U64 (Z.shiftl l r))
      else Result.Err (PartialVMError.new StatusCode.ARITHMETIC_ERROR)
  | IntegerValue.U128 l, IntegerValue.U128 r =>
      if r <? 128
      then Result.Ok (IntegerValue.U128 (Z.shiftl l r))
      else Result.Err (PartialVMError.new StatusCode.ARITHMETIC_ERROR)
  | IntegerValue.U256 l, IntegerValue.U256 r =>
      if r <? 256
      then Result.Ok (IntegerValue.U256 (Z.shiftl l r))
      else Result.Err (PartialVMError.new StatusCode.ARITHMETIC_ERROR)
  | _, _ => Result.Err (PartialVMError.new StatusCode.ARITHMETIC_ERROR)
  end.
\end{lstlisting}
\end{tcolorbox}

\textbf{Explanation}: The \texttt{shl\_checked} function shifts the integer left by a specified number of bits. It ensures that the shift amount does not exceed the bit width of the type (e.g., \texttt{8} for \texttt{U8}). If the shift amount is too large, it raises an arithmetic error.

\begin{tcolorbox}[
  colback=white,     % Background color inside the box (set to white to match the document)
  colframe=white,    % Frame color (set to white to remove it)
  boxrule=0pt,       % No border
  sharp corners,     % Remove rounded corners
  boxsep=0pt,        % Remove padding
  left=0pt,          % Remove left margin
  right=0pt,         % Remove right margin
  top=0pt,           % Remove top margin
  bottom=0pt         % Remove bottom margin
]
\begin{lstlisting}[caption={Casting to \texttt{u64} with \texttt{cast\_u64}},captionpos=b]
Definition cast_u64 (self : IntegerValue.t) : PartialVMResult.t Z :=
  match self with
  | IntegerValue.U8 l => Result.Ok l
  | IntegerValue.U16 l => Result.Ok l
  | IntegerValue.U32 l => Result.Ok l
  | IntegerValue.U64 l => Result.Ok l
  | IntegerValue.U128 l =>
        if l <=? 2^64 - 1
      then Result.Ok l
      else Result.Err (PartialVMError.new StatusCode.ARITHMETIC_ERROR)
  | IntegerValue.U256 l =>
      if l <=? 2^64 - 1
      then Result.Ok l
      else Result.Err (PartialVMError.new StatusCode.ARITHMETIC_ERROR)
  end.
\end{lstlisting}
\end{tcolorbox}




\textbf{Explanation}: The \texttt{cast\_u64} function safely casts values of various integer types to \texttt{U64}. If a \texttt{U256} value exceeds the range of \texttt{U64}, it raises an arithmetic error. Otherwise, it returns the value in a \texttt{Result.Ok}.

%% In a thesis, every section/chapter starts a new page, hence the \clearpage
\clearpage
%% Bibliography/ list of references
%%

\bibliographystyle{unsrt} % Or any other style you prefer (e.g., IEEE, ACM)
\bibliography{references/reference} % Refers to the `references.bib` file

\end{document}